\subsubsection{Scalar constrained operators and conserved coefficients}

For $N=1$, the additive constants of $h_L,h_M$ may be chosen so that
$\mathsf C_L=\mathsf C_M=0$. The matrix operators in
\cref{eq:BA-constrained-operators} are then scalar, with
$\boldsymbol\psi=\psi$ and $\boldsymbol\psi^\vee=\psi^\vee$.
This specialization also exposes two conserved quantities.
Write $I_k=\xi_k\bar\xi_k$, $1\leq k\leq K$. Expanding $\partial^{-1}f=f\partial^{-1}-(\partial f)\partial^{-2}
+O(\partial^{-3})$ and using \cref{eq:general-network-rates} gives
\begin{equation}\label{eq:BA-constrained-H-coefficient}
 \mathcal A_L-\mathcal A_M
 =2\left(\sum_{\ell=1}^{L}I_\ell-\sum_{j=L+1}^{K}I_j\right)\partial^{-1}
      -2H_0\partial^{-2}+O(\partial^{-3}).
\end{equation}
The first two negative-order coefficients of the commutator are
the derivatives of these coefficient differences, up to sign.
Their vanishing recovers conservation of the group power difference
and $H_0$ on the BA realization.
